\section{Breakthrough Target Portfolio}

\subsection{Purpose}

This section identifies the targets where a successful resolution would make later phases disproportionately faster. The goal is to avoid flat exploration. A high-value target is not merely interesting; it should unlock many downstream tests.

\subsection{Method}

The script \texttt{scripts/breakthrough-target-prioritizer.py} combines:

\begin{itemize}
  \item validated probe results;
  \item component contrast groups;
  \item the stem contrast lattice;
  \item sign roles and ready-text counts;
  \item role-boundary controls and blocked phonetic comparisons.
\end{itemize}

It scores targets by downstream unlock value: number of affected probes, corpus coverage, validated grammar constraints, role-boundary risk, and dependency links to later phases.

\subsection{Portfolio Summary}

\begin{table}[htbp]
\centering
\begin{tabular}{lr}
\toprule
\textbf{Metric} & \textbf{Value} \\
\midrule
Breakthrough targets ranked & 9 \\
Dependency edges & 9 \\
Action-plan items & 7 \\
\bottomrule
\end{tabular}
\caption{Breakthrough target summary.}
\label{tab:breakthrough-target-summary}
\end{table}

\subsection{Highest-Leverage Targets}

\begin{table}[htbp]
\centering
\footnotesize
\begin{tabular}{rlp{0.38\textwidth}r}
\toprule
\textbf{Rank} & \textbf{Target} & \textbf{Type} & \textbf{Score} \\
\midrule
1 & \texttt{240} & Grammar keystone & 1.000 \\
2 & \texttt{740::520} & Terminal keystone & 0.828 \\
3 & \texttt{176/482/904/048/061/773} & Stem lattice keystone & 0.788 \\
4 & \texttt{002 X 740} & Frame keystone & 0.737 \\
5 & \texttt{482::904 inside 240+STEM} & Phonetic keystone & 0.726 \\
6 & \texttt{176 with 048/061} & Stem-family keystone & 0.718 \\
7 & \texttt{240::235} & Formula boundary keystone & 0.695 \\
8 & \texttt{032 role split} & Blocker removal & 0.661 \\
\bottomrule
\end{tabular}
\caption{Highest-leverage decipherment targets.}
\label{tab:breakthrough-targets}
\end{table}

\subsection{Interpretation}

The top target is \texttt{240}. If \texttt{240} can be stabilized as a qualifier, title, or classifier-like component, then \texttt{240-482}, \texttt{240-176}, and \texttt{240-740-090} become part of one controlled grammar problem instead of separate curiosities. That makes the stem contrasts easier to test.

The second target is \texttt{740::520}. These signs affect the terminal side of many inscriptions. Solving this contrast would constrain whether the right edge is suffixal, titular, numerical, allographic, or formulaic.

The first phonetic keystone remains \texttt{482::904} inside \texttt{240+STEM}. It is not ranked first because it depends on the grammar environment: if \texttt{240} is not understood well enough, the phonetic test floats without an anchor.

\subsection{Action Plan}

The recommended breakthrough order is:

\begin{enumerate}
  \item Promote all validated \texttt{240} morphology constraints into a grammar model.
  \item Model \texttt{740::520} as competing terminal functions.
  \item Build a stem-lattice consistency test for \texttt{176}, \texttt{482}, \texttt{904}, \texttt{048}, \texttt{061}, and \texttt{773}.
  \item Treat \texttt{002|740} as a formal test environment and score every \texttt{X} filler.
  \item Image-check all \texttt{240-482} and \texttt{240-904} occurrences.
  \item Verify every \texttt{048} and \texttt{061} occurrence in \texttt{002|740}.
  \item Model \texttt{240::235} as a rival opener contrast.
\end{enumerate}

\subsection{Dependency Logic}

The portfolio has a dependency chain:

\begin{center}
\texttt{240 GRAMMAR -> 240+STEM -> 482::904 -> STEM LATTICE -> PHONETIC TESTS}
\end{center}

and a second chain:

\begin{center}
\texttt{740::520 -> 002 X 740 FRAME -> ENTITY-SLOT GRAMMAR}
\end{center}

The main blocker is \texttt{032}. It creates many role-boundary controls. Until its functions are split, it should protect the project from false phonetic comparisons rather than drive new readings.

\subsection{Outputs}

\begin{itemize}
  \item \texttt{outputs/breakthrough\_target\_summary.csv}
  \item \texttt{outputs/breakthrough\_targets.csv}
  \item \texttt{outputs/breakthrough\_dependency\_edges.csv}
  \item \texttt{outputs/breakthrough\_action\_plan.csv}
  \item \texttt{outputs/breakthrough\_target\_portfolio.tex}
\end{itemize}
